%------------------------------------------------------------------------------------
%	CHAPTER 3
%------------------------------------------------------------------------------------
\chapterimage{headConceito.png}
\chapter{EDA e ADS}

\begin{remark}
Nós torturamos os dados até eles confessarem. (Ricado Cappra - Cientista de Dados) 
\end{remark}

A \textbf{Análise e Ciência de Dados} (ADS) é uma área que integra estatística, programação e conhecimento de domínio para transformar grandes volumes de dados em informações úteis e estratégicas. Seu principal objetivo é extrair padrões, tendências e previsões a partir de dados estruturados ou não estruturados, possibilitando a tomada de decisões embasadas. Utilizando técnicas como análise exploratória, aprendizado de máquina e visualização de dados, a \textbf{ADS} atua como ponte entre dados brutos e ações inteligentes.

Mas antes de falarmos dessa prefiro iniciar o assunto com outro acrônimo conhecido como \textbf{EDA}.

\section{Passos da EDA}\index{EDA e ADS}
Acrônimo para \textit{Exploratory Data Analysis} (Análise Exploratória de Dados), desempenha um papel crítico na compreensão do quê? por que? e como? na declaração do problema. \textbf{EDA} é fundamental para entender qualquer conjunto de observações. É aqui podemos obter informações e fazer descobertas. Aqui colocamos o conhecimento para trabalhar.

É a primeira ação realizada na ordem das operações que um Cientista de Dados deve executar ao receber uma nova fonte de observações e a declaração de problema. Tratamos EDA como uma série de técnicas utilizadas como forma de entendermos os diversos aspectos que temos para trabalhar.
\begin{figure}[H]
	\centering
	\includegraphics[width=0.6\textwidth]{cap03/EDA.png}
	\caption{Passos da EDA}
\end{figure}

A preparação dos dados para análise é inevitável, e a maneira como fazemos isso define a sua qualidade. Na prática o que faremos nesse capítulo é compreendermos o que temos a nossa disposição para trabalhar.

Normalmente as observações se dividem em atributos \textbf{Preditores} (Entradas) e \textbf{Alvo} (saída). Uma vez que os localizemos, devemos identificar seu tipo e categoria.

\section{Passo 1 - Entender os Dados}\index{EDA e ADS}
Nesta fase precisamos compreender o que temos a nossa disposição. Começamos o processo com a importação das bibliotecas:
\begin{lstlisting}[]
import pandas as pd
import matplotlib.pyplot as plt
import seaborn as sns
%matplotlib inline
\end{lstlisting}

Vamos trabalhar com três bibliotecas básicas, que como já mencionamos devemos conhecê-las a fundo: \textbf{Pandas} para análise, \textbf{MatPlotLib} e \textbf{SeaBorn} para mostrar em forma gráfica.

Agora precisamos dos dados, para isso usaremos o arquivo \textit{StudentsPerformance.csv}:
\begin{lstlisting}[]
df = pd.read_csv('StudentsPerformance.csv')
\end{lstlisting}

Nessa fase compreendemos melhor o que temos na nossa mão, \textbf{Pandas} é ideal para essa tarefa. Seu funcionamento é como um "Editor de Planilha", dessa forma que devemos encarar essa biblioteca, sua diferença básica é a nomenclatura de como o \textit{DataFrame} (e não Planilha) é visualizado:
\begin{figure}[H]
	\centering
	\includegraphics[width=0.7\textwidth]{cap03/PandasData.png}
	\caption{Visão Pandas}
\end{figure}

Uma coluna aqui é vista como uma \textit{Serie} e \textit{index} é o que mantêm a "cola" das series juntas. Dois comandos são básicos para visualizarmos o \textit{DataFrame}:
\begin{lstlisting}[]
df.head()
\end{lstlisting}

Porém nesse livro não utilizaremos o termo "coluna" e sim "atributo" (consideremos ambos como sinônimos). Mostra as primeiras observações, como parâmetro podemos passar a quantidade. E:
\begin{lstlisting}[]
df.tail()
\end{lstlisting}

Mostra as últimas observações e também como parâmetro podemos passar a quantidade. O que temos até o momento? Sabemos é uma base sobre estudantes e as linhas são: gênero, etnicidade, nível de escolaridade dos pais, forma de alimentação, realizou um teste de preparação do curso, nota de matemática, nota de leitura e nota de escrita.

Então só com esses dois comandos já podemos saber sobre qual assunto iremos tratar: estudantes que realizaram provas e em quais condições. Quantos registros temos a nossa disposição? Ou quais são os nomes dos atributos?
\begin{lstlisting}[]
print("Tamanho: ", df.shape)
print("Nome dos Atributos: ", df.columns)
\end{lstlisting}

As variáveis \textit{shape} e \textit{columns} do \textit{DataFrame} respondem aos questionamentos. De forma mais completa podemos usar:
\begin{lstlisting}[]
df.info()
\end{lstlisting}

Nos mostra inclusive o tipo de cada atributo e se contém ou não elementos nulos. Temos 3 atributos que são do tipo inteiro (\textit{int64}) e podemos analisá-los com o comando:
\begin{lstlisting}[]
df.describe()
\end{lstlisting}

Nos fornece as informações estatísticas básicas como média, desvio padrão, menor valor, máximo, 1º quartil (25\%), 2º quartil ou mediana (50\%), 3º quartil (75\%) e o maior valor. Ou seja, as informações para a montagem de um \textbf{BoxPlot}. Vamos montá-lo para melhor visualizar as observações:
\begin{lstlisting}[]
fig, axes = plt.subplots(1, 3, figsize=(10,4))
axes[0].boxplot(df['math score'])
axes[0].set_title("Matemática")
axes[1].boxplot(df['reading score'])
axes[1].set_title("Leitura")
axes[2].boxplot(df['writing score'])
axes[2].set_title("Escrita")
plt.show()
\end{lstlisting}

E obtemos como resultado:
\begin{figure}[H]
	\centering
	\includegraphics[width=0.55\textwidth]{cap03/BoxPlot.png}
	\caption{BoxPlot das Notas}
\end{figure}

\textbf{Boxplot}\footnote{Diagrama de Caixa se prefere, foi atribuída ao matemático \textbf{John W. Tukey} (1915 –2000), curiosamente algumas literaturas chamam de "\textit{Tukey BoxPlot}", mas se realizar uma pesquisa ninguém sabe ao certo quem criou realmente esse diagrama.} é um gráfico que avalia a distribuição das observações. É formado exatamente com os atributos que mostramos na função \textit{describe()}. Porém suas hastes (inferiores e superiores) se estendem do quartil inferior (ou superior) até o menor valor não inferior (ou superior) ao limite. São calculados da seguinte forma: \vspace{-1em}
\begin{itemize}
	\item Limite inferior: $Q_1 - 1,5 \times (Q_3 - Q_1)$.
	\item Limite superior: $Q_3 + 1,5 \times (Q_3 - Q_1)$.
\end{itemize}

Resumidamente é formado da seguinte maneira:
\begin{figure}[H]
	\centering
	\includegraphics[width=0.6\textwidth]{cap03/estrutBoxplot.png}
	\caption{Estrutura do BoxPlot}
\end{figure}

Esses pontos "discrepantes" podem ocorrer acima ou abaixo dos limites, são chamados de \textit{Outliers}. Não é necessariamente um erro, podemos classificá-lo como uma anomalia curiosa e que merece nossa atenção.

\subsection{Localizar os Outliers}\index{EDA e ADS}
Para achar esses \textit{Outliers} isolamos os três atributos numéricos:
\begin{lstlisting}[]
X = df.iloc[:, 5:8].values
\end{lstlisting}

E criamos um novo \textit{DataFrame} somente com a modificação de nome por um número:
\begin{lstlisting}[]
pd.options.display.float\_format = '{:.1f}'.format
xDF = pd.DataFrame(X)
\end{lstlisting}

Para quê isso serve? A função \textit{describe()} cria um \textit{DataFrame}, podemos percorrê-lo, porém fica muito mais simples se cada atributo for um numeral, pois assim podemos usar um comando \textit{for} para isso:
\begin{lstlisting}[]
z = xDF.describe()
for t in z:
  iqr = z[t][6] - z[t][4]
  extMenor = z[t][4] - (iqr * 1.5)
  extMaior = z[t][6] + (iqr * 1.5)
  print('Para o índice \%d valores devem estar abaixo de \%.2f e acima de \%.2f' \% (t, extMenor, extMaior))
\end{lstlisting}

E obtemos o seguinte resultado: \\
\codigo{Para o índice 0 valores devem estar abaixo de 27.00 e acima de 107.00 \\
Para o índice 1 valores devem estar abaixo de 29.00 e acima de 109.00 \\
Para o índice 2 valores devem estar abaixo de 25.88 e acima de 110.88}

Pelo BoxPlot todos os valores estão abaixo, então para localizá-los:
\begin{lstlisting}[]
matOutliers = (X[:,0] < 27)
df[matOutliers]
\end{lstlisting}

E assim mostramos todas as observações que a nota de matemática (índice 0) é abaixo do valor 27. Proceder de mesmo modo para as notas de leitura (índice 1) e escrita (índice 2) e assim desvendar quais são os \textit{Outliers}.

Podemos também analisar graficamente e visualizar a Distribuição Normal de cada atributo, por exemplo para a nota de Escrita:
\begin{lstlisting}[]
sns.kdeplot(df['writing score'], shade=True)
plt.show()
\end{lstlisting}

Obtemos como resultado:
\begin{figure}[H]
	\centering
	\includegraphics[width=0.5\textwidth]{cap03/NotaEscrita.png}
	\caption{Distribuição das observações para Nota de Escrita}
\end{figure}

E assim verificamos como cada atributo numérico se comporta.

\subsection{Tratar Atributos Categóricos}\index{EDA e ADS}
Sabemos que os primeiros cinco atributos do \textit{Dataframe} são categóricos, porém conforme a função \textit{info()} o tipo delas estás \textit{object}. É interessante mudarmos para o tipo caractere para evitarmos quaisquer problemas futuros.
\begin{lstlisting}[]
df['gender'] = df['gender'].astype(pd.StringDtype())
df['race/ethnicity'] = df['race/ethnicity'].astype(pd.StringDtype())
df['parental level of education'] = df['parental level of education'].astype(pd.StringDtype())
df['lunch'] = df['lunch'].astype(pd.StringDtype())
df['test preparation course'] = df['test preparation course'].astype(pd.StringDtype())
\end{lstlisting}

E ao aplicarmos uma nova chamada a função \textit{info()} vemos que os tipos agora estão corretos. Quantos tipos únicos existem para cada atributo?
\begin{lstlisting}[]
df.nunique()
\end{lstlisting}

Mostra a quantidade de valores não repetidos de cada atributos (inclusive os numéricos). E agora sabemos que temos: 2 gêneros, 5 etnicidades, 6 níveis de escolaridade dos pais, 2 formas de alimentação e 2 tipos para teste de preparação do curso. Mas quem são?
\begin{lstlisting}[]
print("Gênero: ", df['gender'].unique())
print("Etinicidade: ", df['race/ethnicity'].unique())
print("Escolaridade dos Pais: ", df['parental level of education'].unique())
print("Refeição: ", df['lunch'].unique())
print("Realizou Preparatório: ", df['test preparation course'].unique())
\end{lstlisting}

\section{Passo 2 - Limpar os Dados}\index{EDA e ADS}
A limpeza dos dados trata de muitos problemas como informação repetida, valores faltantes (que podem ser descobertos por associação) e inconsistentes. Para esse último tipo o pior caso são os nulos. (In)felizmente essa base está horrível para essa fase e assim pegamos um outro arquivo \textbf{titanic.csv}: 
\begin{lstlisting}[]
df = pd.read_csv('titanic.csv')
df.head()
\end{lstlisting}

Repetimos todo o processo da fase anterior para descobrirmos de que se tratam as observações e descobrimos que são os passageiros (sobreviventes ou não - atributo \textit{Survived} - sendo este atributos alvo) do famoso \textbf{RMS Titanic}, este foi pensado para ser o navio mais luxuoso e seguro de sua época e supostamente "inafundável". Como sabemos em sua viagem inaugural de \textit{Southampton} para \textit{Nova Iorque} afundou no dia 14 de abril de 1912 com mais de 1.500 pessoas a bordo. Porém esta base contém apenas 891 registros.

Ao aplicarmos a função info() percebemos que os atributos \textit{Age} (idade), \textit{Cabin} (número da cabine) e \textit{Embarked} (local de Embarque) possuem valores faltantes. Que valores são esses?
\begin{lstlisting}[]
print(df.isnull().sum())
\end{lstlisting}

Sabemos que faltam: 177 em \textbf{Age}, 687 em \textbf{Cabin} e 2 em \textit{Embarked}. Também podemos mostrar exclusivamente os que faltam, isso é útil para quando temos muitos atributos no modelo:
\begin{lstlisting}[]
null_value_stats = df.isnull().sum(axis=0)
null_value_stats[null_value_stats != 0]
\end{lstlisting}

Ou ainda criar uma função personalizada que retorna um \textit{Dataframe} com a informação mais completa o possível (inclusive com seu percentual):
\begin{lstlisting}[]
def mostrarNulos(data):
  null_sum = data.isnull().sum()
  total = null_sum.sort_values(ascending=False)
  percent = (((null_sum / len(data.index))*100).round(2)).sort_values(ascending=False)
  df_NULL = pd.concat([total, percent], axis=1, keys=['Tot.Nulo', 'Perc.Nulo'])
  df_NULL = df_NULL[(df_NULL.T != 0).any()]
  return df_NULL
\end{lstlisting}

E ao chamá-la:
\begin{lstlisting}[]
df_Age = mostrarNulos(df)
df_Age.head()
\end{lstlisting}

Obtemos como resultado:
\begin{figure}[H]
	\centering
	\includegraphics[width=0.3\textwidth]{cap03/TitanicNulosPerc.png}
	\caption{Nulos e Percetual da Base Titanic}
\end{figure}

Lidar com esses tipos de nulos é complicado pois não temos como consultar e o máximo que podemos fazer é podá-los da nossa base ou atribuir um valor genérico que não afete nosso resultado (como o caso de \textit{Embarked}). Porém \textbf{Número da Cabine} é um dado relevante? Essa é a principal pergunta que nos devemos fazer, por exemplo existe algum modelo preditivo que possa nos dizer que se estivéssemos em determinada cabine no navio sobreviveríamos ou não? Entretanto \textbf{Idade} é um dado relevante (lembra da frase: mulheres e crianças primeiro), então essa é uma característica que pode ser essencial.

Criar uma função com um gráfico para mostrar, por idade, como estão as observações:
\begin{lstlisting}[]
def executarGrafico():
  try:
    sns.distplot([df['Age']])
    plt.show()
  except ValueError as err:
    print(err) 
\end{lstlisting}

Agora a cada vez que chamarmos essa função obtemos como resultado:
\begin{figure}[H]
	\centering
	\includegraphics[width=0.45\textwidth]{cap03/TitanicIdade1.png}
	\caption{Gráfico de Idade do Titanic}
\end{figure}

\begin{note}[Imputação ou retirada de valores] 
	Como tratamos de adicionar ou retirar elementos na base a cada vez devemos ler novamente as observações contidas no arquivo CSV.
\end{note}

Porém em algumas versões da \textit{SeaBorn} este pode apresentar erro devido a presença dos nulos, é ideal que os retiremos do \textit{DataFrame} para evitarmos problemas. Em muitas biografias encontramos algo do tipo: "atribuir um valor (preferencialmente \textit{outlier}) para estes tipos". Tentaremos essa técnica com os seguintes comandos:
\begin{lstlisting}[]
df['Age'].fillna(-25, inplace=True)
executarGrafico()
\end{lstlisting}

Obtemos como resultado:
\begin{figure}[H]
	\centering
	\includegraphics[width=0.45\textwidth]{cap03/TitanicIdade2.png}
	\caption{Gráfico de Idade do Titanic com Outliers}
\end{figure}

Nosso gráfico de idade ganhou uma nova barra, que sabemos com valores não existentes, também podemos atribuir qualquer outro valor como por exemplo a média:
\begin{lstlisting}[]
df['Age'] = df['Age'].fillna(df['Age'].mean())
executarGrafico()
\end{lstlisting}

Ou a mediana (função median()) que resultaria em um gráfico completamente esquisito. Sendo assim vamos cortar esses valores:
\begin{lstlisting}[]
df = df.dropna(axis=0)
executarGrafico()
\end{lstlisting}

E teremos a seguinte situação:
\begin{figure}[H]
	\centering
	\includegraphics[width=0.45\textwidth]{cap03/TitanicIdade3.png}
	\caption{Gráfico de Idade do Titanic sem nulos}
\end{figure}

O que aconteceu? O comando executado eliminou todas as linhas que possuíam valores nulos, e o atributo \textit{Cabin} interferiu e nos deixou, conforme pode ser mostrado com a função \textit{info()}, somente 183 registros no total. Ou seja, o corte que devemos aplicar deve ser cirúrgico e somente no atributo que representa a idade.
\begin{lstlisting}[]
df['Age'] = df['Age'].dropna(axis=0)
executarGrafico()
\end{lstlisting}

O que nos resulta no mesmo gráfico mostrado no início desta e 891 registros. Como citamos, podemos retirar o atributo \textit{Cabin} para que este não interfira mais em futuras análises:
\begin{lstlisting}[]
df = df.drop(['Cabin'], axis=1)
\end{lstlisting}

\begin{note}[Ferramenta para Limpeza dos Dados] 
	Conhece o \textbf{OpenRefine?} é uma ferramenta gratuita dedicada a limpeza e tratamento das observações, baixe uma apostila gratuitamente na minha página do Academia.edu (\url{https://iesbpreve.academia.edu/FernandoAnselmo}).
\end{note}

\section{Passo 3 - Relacionamento entre os Atributos}\index{EDA e ADS}
Vamos retomar nossa base de \textbf{Estudantes} e verificarmos como os atributos se relacionam:
\begin{lstlisting}[]
df = pd.read_csv('StudentsPerformance.csv')
df.corr()
\end{lstlisting}

E temos um valor que corresponde ao grau de relacionamento, um intervalo de -1 a 1, sendo quanto mais próximo do mínimo menor é seu grau de relacionamento. Porém é muito mais fácil de visualizarmos esse resultado com um Mapa de Calor:
\begin{lstlisting}[]
rel = df.corr()
sns.heatmap(rel, xticklabels=rel.columns, yticklabels=rel.columns, annot=True)
plt.show()
\end{lstlisting}

Obtemos como resultado:
\begin{figure}[H]
	\centering
	\includegraphics[width=0.5\textwidth]{cap03/MapaDeCalor.png}
	\caption{Mapa de Calor dos Relacionamentos}
\end{figure}

Vemos que as notas de Escrita e Leitura possuem um forte grau de relacionamento, como se uma fosse a responsável pela outra. Já a de matemática interfere mais na nota de leitura.

Curiosamente se aplicarmos isso na base do \textbf{Titanic} vemos que os atributos mais importantes para \textit{Fare} (sobreviveu) que é nosso alvo são: \textit{Fare} que é o valor pago pela passagem e \textit{Parch} que se refere a quantidade de pais. Ou seja, os mais ricos e se a criança tinha ou não os pais a bordo de modo a colocá-las no bote salva vidas. 

Outra maneira de visualizarmos, também de forma gráfica, é através da dispersão de valores:
\begin{lstlisting}[]
sns.pairplot(df)
plt.show()
\end{lstlisting}

Obtemos como resultado:
\begin{figure}[H]
	\centering
	\includegraphics[width=0.7\textwidth]{cap03/dispersao.png}
	\caption{Dispersão Associada}
\end{figure}

Quanto mais juntos aparecem os pontos mais relacionadas estão. Podemos isolar as notas de Escrita e Leitura em um único gráfico, por exemplo:
\begin{lstlisting}[]
sns.regplot(x='writing score', y='reading score', data=df)
plt.show()
\end{lstlisting}

Esta função executa um ajuste e plotagem simples com base no modelo de Regressão Linear. E obtemos como resultado:
\begin{figure}[H]
	\centering
	\includegraphics[width=0.55\textwidth]{cap03/NotasRegress.png}
	\caption{Notas de Leitura e Escrita}
\end{figure}

Porém o mais interessante é colorir os pontos de forma diferente com base em um atributo categórico que pode ser uma causa (para uma nota alta ou baixa), por exemplo o quanto a alimentação interferiu na nota:
\begin{lstlisting}[]
sns.lmplot(x='writing score', y='reading score', hue='lunch', data=df)
plt.show()
\end{lstlisting}

A função \textit{lmplot()} combina \textit{regplot()} com a classe \textbf{FacetGrid}. Esta auxilia visualizar a distribuição de um determinado atributos, bem como o relacionamento entre os vários separadamente dentro de subconjuntos das observações. Obtemos como resultado:
\begin{figure}[H]
	\centering
	\includegraphics[width=0.5\textwidth]{cap03/NotasLmplot.png}
	\caption{Nota associada a Alimentação - RegPlot}
\end{figure}

Uma melhor forma de visualizar é usar a função \textit{relplot()} que fornece acesso a várias funções diferentes no nível de eixos que mostram o relacionamento entre dois atributos com mapeamentos semânticos de subconjuntos:
\begin{lstlisting}[]
sns.relplot(x='writing score', y='reading score', hue='lunch', data=df)
plt.show()
\end{lstlisting}

Obtemos como resultado:
\begin{figure}[H]
	\centering
	\includegraphics[width=0.55\textwidth]{cap03/NotasAlimentacao.png}
	\caption{Nota associada a Alimentação - RelPlot}
\end{figure}

Ou seja, podemos responder várias perguntas apenas com a verificação do relacionamento entre os atributos. Como forma de fixar o conhecimento procure realizar o mesmo teste com outros atributos categóricos e descobrir como se comportam em relação a nota, se existe ou não interferência.

\section{ADS}\index{EDA e ADS}
\textit{Analysis and Data Science} (Análise e Ciência de Dados) trata de uma área multidisciplinar que combina técnicas de estatística, matemática, programação e conhecimento do domínio para extrair \textit{insights} e gerar valor a partir de dados. Seu objetivo é transformar dados brutos em valor, tais como, previsões inteligentes e decisões orientadas por evidências.

Com o crescimento exponencial da geração de dados por empresas, redes sociais, dispositivos móveis e sensores, tornou-se essencial utilizar abordagens analíticas para compreender padrões, detectar anomalias, prever comportamentos e otimizar processos.\textbf{}

Na prática de ADS geralmente segue etapas conhecidas como pipeline de ciência de dados: \begin{itemize}[]
	\item \textbf{Coleta de dados} - Captura de dados de fontes diversas (bancos de dados, APIs, arquivos, web scraping).
	\item \textbf{Limpeza e pré-processamento} - Tratamento de dados ausentes, inconsistentes ou ruidosos.
	\item \textbf{Análise exploratória (EDA)} - Visualizações e estatísticas para entender padrões, correlações e distribuições.
	\item \textbf{Engenharia de atributos (\textit{feature engineering})} - Transformação de variáveis para melhor desempenho dos modelos.
	\item \textbf{Modelagem estatística e de machine learning} - Criação de modelos preditivos ou descritivos.
	\item \textbf{Validação e avaliação} - Medição do desempenho usando métricas como acurácia, precisão, recall ou RMSE.
	\item \textbf{Deploy e monitoramento} - Aplicação do modelo em produção e acompanhamento de seu desempenho.
\end{itemize}

\section{PDP e ICE}\index{EDA e ADS}
No contexto da \textbf{ADS}, os gráficos \textbf{PDP} (\textit{Partial Dependence Plot}) e \textbf{ICE} (\textbf{Individual Conditional Expectation}) são ferramentas úteis para interpretar modelos de Machine Learning, especialmente os mais complexos (como Árvores de Decisão, Random Forests ou Gradient Boosting). Ambos ajudam a entender como uma variável de entrada influencia a predição de um modelo.

\textbf{PDP - Partial Dependence Plot (Gráfico de Dependência Parcial)} \\
Mostra a média da predição do modelo ao variar uma variável de interesse (ou mais), mantendo as outras constantes. Responde à pergunta: Em média, como a predição muda quando alteramos o valor de uma variável específica?

Como funciona? \vspace{-1em}
\begin{enumerate}[]
	\item Para cada valor da variável de interesse: \vspace{-1em}
	\begin{itemize}[]
		\item Substitui-se esse valor em todas as amostras do conjunto de dados.
		\item O modelo faz a predição para cada uma dessas amostras modificadas.
		\item Calcula-se a média das predições.
	\end{itemize}
	\item Isso é feito para vários valores da variável (ex: de 0 a 10), e o resultado é um gráfico com uma curva.
\end{enumerate}

\textbf{CE – Individual Conditional Expectation (Expectativa Condicional Individual)} \\
É uma generalização do \textbf{PDP}. Ao invés de mostrar a média, mostra uma linha para cada observação individual. Ele responde à pergunta: Como a predição muda para um indivíduo específico quando mudamos uma variável?

Como funciona? \vspace{-1em}
\begin{enumerate}[]
	\item Para cada amostra individual:
	\begin{itemize}[]
		\item Varia-se uma variável específica mantendo as outras constantes.
		\item Registra-se a predição do modelo para cada valor variado.
	\end{itemize}
	\item Plota-se uma curva por observação → o resultado é um conjunto de curvas (um para cada indivíduo).
\end{enumerate}

Imaginemos um modelo que prediz o preço de uma casa com base em tamanho ($m^2$): \vspace{-1em}
\begin{itemize}[]
	\item \textbf{PDP}: Mostra que, em média, quanto maior a casa, maior o preço.
	\item \textbf{ICE}: Mostra que, para algumas casas, o aumento de tamanho aumenta o preço rapidamente, para outras, quase não muda — indicando efeitos individuais distintos.
\end{itemize}

\begin{lstlisting}[]
import matplotlib.pyplot as plt
import pandas as pd
from sklearn.datasets import fetch_california_housing
from sklearn.ensemble import RandomForestRegressor
from sklearn.inspection import plot_partial_dependence, PartialDependenceDisplay
from sklearn.inspection import partial_dependence
from sklearn.model_selection import train_test_split
from sklearn.preprocessing import StandardScaler

# 1. Carregar o dataset
data = fetch_california_housing()
X = pd.DataFrame(data.data, columns=data.feature_names)
y = data.target

# 2. Dividir em treino e teste
X_train, X_test, y_train, y_test = train_test_split(X, y, random_state=42)

# 3. Treinar o modelo
model = RandomForestRegressor(n_estimators=100, random_state=42)
model.fit(X_train, y_train)

# 4. Plotar PDP + ICE para uma variável
features = ['MedInc']

fig, ax = plt.subplots(figsize=(10, 6))

display = PartialDependenceDisplay.from_estimator(
 model,
 X_train,
 features,
 kind="both",
 subsample=50,
 n_jobs=1,
 grid_resolution=100,
 ax=ax
)

plt.title("PDP + ICE para a variável MedInc")
plt.grid(True)
plt.tight_layout()
plt.show()
\end{lstlisting}

Acabei de expor o grande segredo de um Modelo de \textbf{Machine Learning}, que nada mais é do que, obter os dados, separar em treino e teste, executar o modelo. Mas veremos isso mais a fundo nos próximos capítulos.

Basicamente o que nos interessa é o \textbf{Passo 4} no qual, neste exemplo, usamos a variável \textbf{MedInc} (renda média da região), mas se desejar podemos trocar para \textbf{AveRooms}, \textbf{AveOccup} ou \textbf{HouseAge}. Já o método "\textit{from\_estimator}", o parâmetro \textbf{kind} pode ser usado com os valores: '\textit{average}' para \textbf{PDP}, '\textit{individual}' para \textbf{ICE} ou '\textit{both}' para os dois. O parâmetro \textit{subsample} permite limitar uma quantidade de indivíduos no ICE para visualização (neste caso estabelecido 50).\textbf{}

Obtemos como resultado:
\begin{figure}[H]
	\centering
	\includegraphics[width=0.55\textwidth]{cap03/PDPeICE.png}
	\caption{PDP e ICE para variável MedInc}
\end{figure}

Na prática o \textbf{PDP} é usado para Modelos de \textbf{Classificação} enquanto que o \textbf{ICE} para Mdelos de \textbf{Regressão}.

A curva azul escura (grossa) é a média das previsões (PDP). Já as muitas linhas finas e transparentes cada uma representa uma instância (ICE). O eixo X representa o valor da variável MedInc (renda média), enquanto que o eixo Y mostra como a predição (\textit{target}) muda conforme MedInc.
	
\section{Conclusão}\index{EDA e ADS}
Um profissional de \textbf{ADS} desenvolve:
\begin{itemize}[nolistsep]
	\item Pensamento analítico e crítico
	\item Capacidade de resolver problemas baseados em dados
	\item Habilidade de comunicar insights com clareza (data storytelling)
	\item Conhecimentos em estatística, probabilidade e algoritmos de aprendizado de máquina
\end{itemize}

Mantemos em mente que \textbf{EDA} e um aspecto central da \textit{Data Science}, que às vezes é esquecido. O primeiro passo de qualquer ação que tomemos é conhecer as observações: entendê-las e familiarizar-se. Quais são as respostas que estamos tentando obter? Quais são os atributos e o que significam? Como é a aparência de uma perspectiva estatística? As observações estão formatadas corretamente? Possuem valores ausentes? duplicados? E quanto aos \textit{outliers}? Conhecemos eles? Ou seja, devemos responder a esses questionamentos.

É necessário muito trabalho de preparação, pois no mundo real dados raramente são limpos e homogêneos. Costumamos dizer que 80\% do tempo valioso em um Cientista de Dados é utilizado com a localização, limpeza e organização das observações. Os 20\% restantes são destinados a realizar as análises.

O campo de \textbf{ADS} é essencial em um mundo cada vez mais orientado por dados, sendo aplicado em setores como saúde, finanças, varejo, indústria e governo. Com o avanço da tecnologia e o crescimento massivo na geração de dados, a demanda por profissionais capazes de analisar e interpretar essas informações tem aumentado exponencialmente. A ciência de dados não apenas responde a “o que aconteceu”, mas também busca responder “por que aconteceu” e “o que pode acontecer”, contribuindo para a inovação, otimização de processos e ganho de vantagem competitiva.

Agora estamos prontos para começarmos a explorar diversas observações com a utilização dos modelos.

\clearpage